%% Sample Exam 5 -- Medical Image Analysis (ETH Zurich, 2026)
%%
%% Written in the style of Example_Exam_with_Answers.pdf.
%%
%% Compile the version WITH solutions:      pdflatex "\def\withanswers{}%% Sample Exam 5 -- Medical Image Analysis (ETH Zurich, 2026)
%%
%% Written in the style of Example_Exam_with_Answers.pdf.
%%
%% Compile the version WITH solutions:      pdflatex "\def\withanswers{}%% Sample Exam 5 -- Medical Image Analysis (ETH Zurich, 2026)
%%
%% Written in the style of Example_Exam_with_Answers.pdf.
%%
%% Compile the version WITH solutions:      pdflatex "\def\withanswers{}%% Sample Exam 5 -- Medical Image Analysis (ETH Zurich, 2026)
%%
%% Written in the style of Example_Exam_with_Answers.pdf.
%%
%% Compile the version WITH solutions:      pdflatex "\def\withanswers{}\input{Sample_Exam_5.tex}"
%% Compile the blank version (no answers):  pdflatex Sample_Exam_5.tex
%%
\ifx\withanswers\undefined
  \documentclass[11pt]{exam}
\else
  \documentclass[11pt,answers]{exam}
\fi

\usepackage[margin=2.4cm]{geometry}
\usepackage{amsmath,amssymb}
\usepackage{array}
\usepackage[T1]{fontenc}

% ---------------------------------------------------------------- page style
\pagestyle{foot}
\firstpageheader{}{}{}
\firstpagefooter{}{}{}
\runningheader{}{}{}
\runningfooter{}{}{Page \thepage}

% -------------------------------------------------------------- multiplechoice
% Reproduce the "(A) (B) (C) ..." labelling of the original exam.
\renewcommand{\choicelabel}{$\bigcirc$~(\Alph{choice})}

% ---------------------------------------------------------- true/false macros
% \TFT -> the statement is TRUE, \TFF -> the statement is FALSE.
% In the answer version the correct word is set in capitals and bold.
\newcommand{\TFT}{\ifprintanswers\textbf{TRUE}\hspace{1.1em}False\else True\hspace{1.1em}False\fi}
\newcommand{\TFF}{True\hspace{1.1em}\ifprintanswers\textbf{FALSE}\else False\fi}

% ------------------------------------------------------------- writing space
% Reserves room to write in the blank version; collapses in the answer version.
\newcommand{\answerspace}[1]{\ifprintanswers\else\vspace{#1}\fi}

\begin{document}

% ------------------------------------------------------------------ title page
\begin{center}
  {\LARGE\bfseries Sample Exam}\\[2mm]
  {\Large Medical Image Analysis}\\[6mm]
\end{center}

\vspace{4mm}
\noindent Student name:\hrulefill

\vspace{6mm}
\noindent Immatriculation number:\hrulefill

\vspace{10mm}

\begin{center}
  {\large\bfseries Exam Questions}
\end{center}

\noindent\textbf{Note:} All the questions have equal weight.

\vspace{4mm}

\begin{questions}

% =============================================================== Question 1 ===
\question
A measurement taken in an image corresponds to a real-world quantity only through the voxel size.
An anatomical structure is segmented in two scans of the same patient, acquired one year apart on
different scanners.

\begin{parts}

  \part In the \emph{baseline} scan the voxel size is $0.8 \times 0.8 \times 3.0$~mm$^3$ and the
  segmentation contains $4\,500$ voxels. What is the volume of the structure, in mm$^3$ and in
  cm$^3$?

  \begin{solution}
    Answer: $8\,640$~mm$^3$ $= 8.64$~cm$^3$.

    One voxel occupies $0.8 \times 0.8 \times 3.0 = 1.92$~mm$^3$, so
    $4\,500 \times 1.92 = 8\,640$~mm$^3$. There are $1\,000$~mm$^3$ in a cm$^3$.
  \end{solution}
\answerspace{2.4cm}

  \part In the \emph{follow-up} scan the voxel size is $0.5 \times 0.5 \times 2.0$~mm$^3$ and the
  segmentation contains $16\,000$ voxels. What is the volume now, and has the structure grown?

  \begin{solution}
    Answer: $8\,000$~mm$^3$ $= 8.00$~cm$^3$. The structure has \emph{not} grown --- it has shrunk
    slightly.

    A voxel now occupies $0.5 \times 0.5 \times 2.0 = 0.5$~mm$^3$, so
    $16\,000 \times 0.5 = 8\,000$~mm$^3$. Compared with the baseline this is a change of
    $(8\,000 - 8\,640)/8\,640 = -7.4\%$.
  \end{solution}
\answerspace{2.8cm}

  \part What would you have concluded if you had compared the raw voxel counts instead, and what
  does this tell you about comparing measurements across acquisitions?

  \begin{solution}
    Answer: you would have concluded that the structure grew by a factor of
    $16\,000 / 4\,500 \approx 3.6$, i.e.\ by more than $250\%$ --- the opposite of the truth, and
    wrong by a wide margin.

    A voxel count is not a measurement of anything physical. It becomes one only after being
    multiplied by the voxel volume, and the voxel volume differs between the two acquisitions here
    by almost a factor of four. Whenever measurements from different images are compared ---
    longitudinal studies, population statistics, normative distributions --- the real-world voxel
    size has to be taken into account, either by converting to physical units as here, or by
    resampling the images to a common voxel size first.
  \end{solution}
\answerspace{4.2cm}

\end{parts}

% =============================================================== Question 2 ===
\question
Deep learning libraries provide a layer called ``convolution'', but the operation they implement is
usually the closely related cross-correlation. For an image $I$ and a kernel $k$, which of the two
expressions below is the \emph{convolution} $I * k$? Please circle the correct choice.

\vspace{2mm}
\begin{center}
  $\displaystyle \sum_m \sum_n I(x-m,\, y-n)\, k(m,n)$
  \hspace{2.0cm}
  $\displaystyle \sum_m \sum_n I(x+m,\, y+n)\, k(m,n)$
\end{center}
\vspace{2mm}

\begin{solution}
  Answer: Left one.

  Convolution reverses one of the two functions before shifting it, which is what the minus signs
  encode; the right-hand expression has no flip and is the cross-correlation. The only thing the
  flip buys is commutativity, $I * k = k * I$, which is of little practical use in a neural network
  --- and since the kernel weights are learned anyway, a flipped kernel is just as learnable as an
  unflipped one. That is why most libraries implement cross-correlation and call it convolution,
  and why the distinction almost never matters in practice. It does matter when you take a kernel
  from a signal-processing textbook and expect it to behave identically.
\end{solution}
\answerspace{1.5cm}

% =============================================================== Question 3 ===
\question
A great many methods in this course have the same shape: a data term measuring agreement with the
observation, plus a regularisation term expressing which solutions are plausible,

\vspace{1mm}
\begin{center}
  $\displaystyle \min_\theta\ D(\text{data};\theta) \; + \; \lambda\, R(\theta).$
\end{center}
\vspace{1mm}

\noindent For each of the three methods below, please state what plays the role of the data term
and what plays the role of the regulariser.

\begin{parts}
  \part Total variation denoising.
  \begin{solution}
    Answer: data term $D = \|I - J\|_2^2$, the squared difference between the recovered image and
    the noisy observation. Regulariser $R = \|I\|_{TV} = \int \|\nabla I\|\,dx$, the total
    variation, which prefers images with minimal variation and therefore a piecewise-constant
    tendency.
  \end{solution}
\answerspace{2.6cm}

  \part Deformable intensity-based registration with linear-elastic regularisation.
  \begin{solution}
    Answer: data term $D = d(I,\, T_\theta \circ J)$, an intensity similarity metric such as the
    sum of squared differences between the target and the warped moving image. Regulariser
    $R = r(T_\theta^{-1})$, the elastic strain energy of the displacement field
    $u$, built from the Cauchy strain tensor $V(u) = (\nabla u + \nabla u^T)/2$, which prefers
    small, smooth deformations.
  \end{solution}
\answerspace{2.8cm}

  \part MAP segmentation with a Potts / Ising Markov random field prior.
  \begin{solution}
    Answer: data term is the unary term
    $\sum_x (I(x)-\mu_{c(x)})^T \Sigma^{-1}_{c(x)} (I(x)-\mu_{c(x)})$, i.e.\ the negative
    log-likelihood of the intensity at each pixel under the class it was assigned. Regulariser is
    the pairwise term $\lambda \sum_x \sum_{y \in G(x)} \delta(c(x) \neq c(y))$, which charges a
    fixed penalty for every neighbouring pair of pixels carrying different labels and therefore
    prefers spatially consistent labellings.
  \end{solution}
\answerspace{3.0cm}
\end{parts}

% =============================================================== Question 4 ===
\question
Please indicate whether the statement is true or false by circling the correct choice in each of
the following.

\begin{parts}
  \part \TFT\ \ Convolution is translation equivariant but not translation invariant.

  \part \TFF\ \ A ``same'' (padded) convolution with a $k \times k$ kernel and stride 1 reduces
  each spatial dimension by $k-1$.

  \part \TFF\ \ Fully connected layers are naturally translation invariant, which is why they are
  placed at the end of a classification network.

  \part \TFT\ \ For a recognition task two images differing only by a translation should give
  identical outputs, whereas for segmentation they should give the same output at a translated
  location.
\end{parts}

\begin{solution}
  (a) TRUE --- and the distinction is worth being precise about. Equivariance means
  $f(T(x)) = T(f(x))$: translate the input and the feature map translates with it. Invariance would
  mean $f(T(x)) = f(x)$, the output not changing at all, and convolution does \emph{not} give that.
  Pooling is what contributes a partial, local invariance.

  (b) FALSE --- that is the \emph{valid} convolution, $N_l = N_{l-1} - k_1 + 1$. The whole purpose
  of ``same'' convolution is to pad the boundary so the output has the same size as the input,
  $N_l = N_{l-1}$.

  (c) FALSE --- exactly backwards. Translation invariance is \emph{not} native to fully connected
  layers: translating the object produces completely different hidden activations, because every
  input position has its own separate weight. One can teach an FC network some invariance by
  augmenting the training data with random translations, but the lecture's verdict on that is
  ``not the most elegant way''.

  (d) TRUE --- and this is why the two properties are both wanted, in different places.
  Recognition wants invariance: a tumour is a tumour wherever it sits. Segmentation and detection
  want equivariance: move the structure and the output label map should move with it.
\end{solution}

% =============================================================== Question 5 ===
\question
A multi-head self-attention block has input dimension $D = 512$ and $H = 8$ heads, with per-head
dimensions $D_K = D_Q = D_V = 64$. Counting the projection matrices $W_{K_h}, W_{Q_h}, W_{V_h}$ of
every head together with the output projection $W_0$, and ignoring all biases, how many parameters
does the block have?

\begin{solution}
  Answer: $\mathbf{1\,048\,576}$ parameters.

  Per head, each of $W_{K_h}, W_{Q_h}, W_{V_h}$ has shape $D \times 64 = 512 \times 64 = 32\,768$
  entries, so a head contributes $3 \times 32\,768 = 98\,304$. Over 8 heads:
  $8 \times 98\,304 = 786\,432$.

  The heads are concatenated, giving $H \cdot D_V = 8 \times 64 = 512$ features, which the output
  projection maps back to $D$: $W_0 \in \mathbb{R}^{512 \times 512} = 262\,144$ entries.

  Total: $786\,432 + 262\,144 = 1\,048\,576$. Note that with $H \cdot D_V = D$, which is the usual
  choice, this is exactly $4D^2$ --- the cost of the block is quadratic in the model width and does
  not depend on the number of heads.
\end{solution}
\answerspace{4.0cm}

% =============================================================== Question 6 ===
\question
Once a Markov random field prior is added, segmentation can no longer be solved by deciding each
pixel independently. Please explain briefly why exact inference becomes hard, and describe one
approximate method presented in the lecture. You do not need to write down the equations.

\begin{solution}
  Answer: \emph{Why it is hard.} Without the MRF the posterior factorises over pixels, so each
  pixel's label can be chosen independently and the answer is found in one pass. The pairwise term
  couples neighbouring labels, so the posterior no longer factorises: the best label for a pixel
  depends on its neighbours' labels, which depend on theirs, and so on across the whole image.
  Computing the normalisation constant would require summing over every possible labelling of the
  image, of which there are exponentially many, and exact energy minimisation with categorical
  labels is NP-hard in 2D and above.

  \emph{An approximate method --- Gibbs sampling.} Sampling from the full posterior is infeasible,
  but sampling from the conditional distribution of a \emph{single} voxel given all the others is
  easy, because that conditional only involves the voxel's likelihood and its immediate
  neighbours, and normalising it means summing over the few possible labels of one voxel. So:
  initialise the labelling (for instance from the likelihood alone), then repeatedly visit each
  voxel and re-sample its label from that conditional. It is simple to implement and effective, but
  it is stochastic --- every run gives a different result --- the outcome depends on the
  initialisation and the parameters, and convergence can be slow.

  Iterated conditional modes and graph cuts are the other standard alternatives.
\end{solution}
\answerspace{7.5cm}

% =============================================================== Question 7 ===
\question
Which of the following statements on the differences between Active Shape Models (ASM) and Active
Appearance Models (AAM) is \textbf{inaccurate}?

\begin{choices}
  \choice ASM searches \emph{around} the current position, sampling along the surface normals of
  the model points, while AAM evaluates its fit \emph{at} the current position.
  \choice ASM minimises the distance between the model points and the corresponding points found
  in the image, whereas AAM minimises the difference between a synthesised model image and the
  target image.
  \choice AAM usually has a larger capture range than ASM, because the additional texture
  information gives it more to lock onto when the initialisation is poor.
  \choice AAM builds its texture model by warping every training image to the mean shape before
  sampling and normalising the intensities.
\end{choices}

\begin{solution}
  Answer: C.

  It is \textbf{ASM} that has the larger capture range, and (A) explains why: by sampling along the
  normals it actively looks a certain distance away from where the model currently sits, so it can
  find a boundary that is not underneath the model yet. AAM only compares its synthesised
  appearance with the image at the current position, so if the initialisation is too far off there
  is nothing to pull it in the right direction. ASM is also the faster of the two and locates
  feature points more accurately; what AAM buys in exchange is a better match to the image texture.
\end{solution}

% =============================================================== Question 8 ===
\question
Which of the following statements on non-linear transformation models is \textbf{incorrect}?

\begin{choices}
  \choice A naive parameterisation with one independent displacement vector per voxel can generate
  transformations that are physically implausible and even change the topology of the anatomy.
  \choice In a flow model each point follows the ordinary differential equation
  $dx/dt = v(x,t)$, and if the velocity field is smooth the resulting transformation is guaranteed
  to be a diffeomorphism.
  \choice A stationary velocity field has more parameters than a time-dependent one, which is why
  the scaling-and-squaring scheme is needed to make it tractable.
  \choice The diffusion model requires the displacement field to satisfy the Poisson equation,
  $\Delta u(x') + F(x') = 0$.
\end{choices}

\begin{solution}
  Answer: C.

  A stationary velocity field has \emph{fewer} parameters, not more: dropping the time dependence
  means storing one velocity vector per voxel instead of one per voxel \emph{per time point}. And
  scaling-and-squaring is not a workaround for a cost problem --- it is an efficient integration
  scheme that becomes available precisely \emph{because} the field is stationary. Both halves of
  the statement are inverted.

  The value of (B) is worth remembering separately: a diffeomorphism is smooth, invertible and has
  a smooth inverse, so a transformation built this way physically cannot tear or fold tissue. This
  is the guarantee that the naive parameterisation in (A) lacks, and it is the foundation of LDDMM.
\end{solution}

% =============================================================== Question 9 ===
\question
Which of the following statements on intensity normalisation is \textbf{incorrect}?

\begin{choices}
  \choice Min-max normalisation and matching the mean and standard deviation are both linear
  mappings of the intensities, and neither can change the contrast of an image.
  \choice Nyul's method matches the \emph{locations} of landmarks of the histograms, which aligns
  the distributions without making the two histograms identical.
  \choice Histogram matching exploits the spatial correlation between neighbouring pixels, which is
  what makes it robust to differences in the size of the imaged object.
  \choice Intensity variation is a particular problem for MRI because it has no absolute intensity
  scale, unlike CT where intensities are expressed in Hounsfield units.
\end{choices}

\begin{solution}
  Answer: C.

  Both halves are wrong. Histogram matching knows nothing at all about image content: it treats
  every pixel as an independent sample and uses \emph{no} spatial information or correlation
  between neighbours --- which is why matching image \emph{gradients} has been proposed as a
  remedy. And differences in object size are a genuine problem for it rather than something it is
  robust to: if one image contains much more of a given tissue than the other, a perfect histogram
  match is not even desirable, and mode matching is preferred.

  The other caveat worth knowing alongside this one is pathology: a large tumour changes the
  histogram, so matching a pathological image to a healthy one is hard, and robust statistics or
  outlier detection help.
\end{solution}

% ============================================================== Question 10 ===
\question
In patch-based label fusion the label at a target voxel is decided by
$L(x) = \arg\max_l \sum_i w_i\, \delta(\text{label}_i = l) \,/\, \sum_i w_i$, where the weights
$w_i$ come from the similarity between the patch around the target voxel and patches taken from
the atlas library. Which of the following statements about this scheme is correct?

\begin{choices}
  \choice The weights are computed from the spatial distance between the target voxel and the atlas
  voxel, not from the image content.
  \choice Only a coarse rigid or affine registration of the atlases is required, because
  correspondence is re-established locally by comparing patches inside a search window around each
  voxel.
  \choice The scheme requires exactly one atlas; with several atlases the weights no longer sum to
  one and the decision becomes ill-defined.
  \choice Because the weighting is borrowed from non-local means, the scheme can only be applied to
  denoising and not to the propagation of labels.
  \choice The search window has to cover the whole image, since restricting it would bias the
  weights.
\end{choices}

\begin{solution}
  Answer: B.

  This is the point of patch-based methods. In classical multi-atlas segmentation the
  correspondence between the atlas and the patient is whatever the registration says it is, so the
  result is only as good as a non-linear registration --- which is expensive and can fail. Patch
  matching replaces that: after a coarse alignment, each target voxel looks around a local search
  window for atlas patches that actually resemble its own neighbourhood, and takes their labels
  weighted by that resemblance. Correspondence is decided by image content at the point of use
  rather than by a global transformation computed in advance.

  (E) is the reverse of the truth: restricting the search window is standard practice, for cost, and
  it also encodes the sensible prior that the correct correspondence is somewhere nearby.
\end{solution}

% ============================================================== Question 11 ===
\question
A U-Net for segmentation takes $128 \times 128$ inputs with one channel. All convolutions are
$3 \times 3$ with ``same'' padding; each encoder level ends with $2 \times 2$ max pooling of stride
2, and each decoder level begins by upsampling by a factor of 2. The encoder levels produce 32, 64
and 128 channels respectively, and the bottleneck produces 256 channels. Each skip connection
concatenates the encoder features onto the upsampled decoder features at the matching level. Circle
True or False:

\begin{parts}
  \part \TFT\ \ The feature maps at the bottleneck have a spatial size of $16 \times 16$.

  \part \TFF\ \ Because all convolutions use ``same'' padding, no spatial information is lost in
  the encoder, and the skip connections are therefore redundant.

  \part \TFT\ \ After upsampling the bottleneck to $32 \times 32$ and concatenating the skip
  connection from the matching encoder level, the first decoder convolution at that level receives
  $256 + 128 = 384$ input channels.

  \part \TFF\ \ A $3 \times 3$ convolution mapping those 384 input channels to 128 output channels
  has $3 \cdot 3 \cdot 384 = 3\,456$ weights, excluding biases.
\end{parts}

\begin{solution}
  (a) TRUE --- three poolings, each halving: $128 \rightarrow 64 \rightarrow 32 \rightarrow 16$.
  The ``same'' convolutions do not change the size, so only the pooling does.

  (b) FALSE --- and this is the point of the architecture. It is the \emph{pooling}, not the
  convolutions, that destroys spatial detail: by the bottleneck the feature maps are $16 \times 16$
  and the fine boundary information is gone. The skip connections exist to carry the
  high-resolution features from the encoder across to the decoder, so the network can combine the
  large receptive field of the deep layers with the spatial precision of the shallow ones.

  (c) TRUE --- concatenation adds the channel counts: 256 upsampled from the bottleneck plus the
  128 arriving through the skip connection.

  (d) FALSE --- the count leaves out the output channels. A convolution kernel connecting every
  input channel to every output channel has $k_1 \cdot k_2 \cdot C_{in} \cdot C_{out}$ weights, so
  the correct number is $3 \cdot 3 \cdot 384 \cdot 128 = 442\,368$ --- 128 times the stated figure.
  Note also that this is where much of a U-Net's parameter budget goes: the concatenation makes the
  decoder convolutions wide.
\end{solution}

% ============================================================== Question 12 ===
\question
Which of the following statements on interpretability and its uses are correct (multiple answers
are possible)?

\begin{choices}
  \choice Saliency maps change over the course of training, and can be used as a kind of
  fingerprint of how a model responds to its inputs.
  \choice In interpretability-guided active learning, a sample is treated as informative when the
  model appears to struggle with it; easy samples tend to produce clean, focused saliency.
  \choice A model-agnostic method such as LIME explains a prediction by fitting a simpler,
  interpretable model that approximates the classifier in a neighbourhood of the input, weighted by
  a proximity measure.
  \choice Class-prototype maximisation returns the image from the training set that most strongly
  activates a given class.
  \choice Grad-CAM can be applied without any access to the internals of the model, since it only
  requires the class scores.
\end{choices}

\begin{solution}
  Answer: (A), (B) and (C) are correct.

  (D) misdescribes the method. Class-prototype maximisation performs gradient \emph{ascent on the
  input} to synthesise a pattern that maximises a class activation; the result is not a training
  image and indeed is not necessarily a realistic image at all. Regularising it with a generative
  prior is what makes the synthesised inputs look natural.

  (E) puts Grad-CAM on the wrong side of the taxonomy. It is a \emph{gradient}-based method, and it
  needs the gradients of the class score with respect to the convolutional feature maps --- so it
  requires access to the internals and is a white-box method. LIME is the model-agnostic,
  black-box one.
\end{solution}

\end{questions}

\end{document}
"
%% Compile the blank version (no answers):  pdflatex Sample_Exam_5.tex
%%
\ifx\withanswers\undefined
  \documentclass[11pt]{exam}
\else
  \documentclass[11pt,answers]{exam}
\fi

\usepackage[margin=2.4cm]{geometry}
\usepackage{amsmath,amssymb}
\usepackage{array}
\usepackage[T1]{fontenc}

% ---------------------------------------------------------------- page style
\pagestyle{foot}
\firstpageheader{}{}{}
\firstpagefooter{}{}{}
\runningheader{}{}{}
\runningfooter{}{}{Page \thepage}

% -------------------------------------------------------------- multiplechoice
% Reproduce the "(A) (B) (C) ..." labelling of the original exam.
\renewcommand{\choicelabel}{$\bigcirc$~(\Alph{choice})}

% ---------------------------------------------------------- true/false macros
% \TFT -> the statement is TRUE, \TFF -> the statement is FALSE.
% In the answer version the correct word is set in capitals and bold.
\newcommand{\TFT}{\ifprintanswers\textbf{TRUE}\hspace{1.1em}False\else True\hspace{1.1em}False\fi}
\newcommand{\TFF}{True\hspace{1.1em}\ifprintanswers\textbf{FALSE}\else False\fi}

% ------------------------------------------------------------- writing space
% Reserves room to write in the blank version; collapses in the answer version.
\newcommand{\answerspace}[1]{\ifprintanswers\else\vspace{#1}\fi}

\begin{document}

% ------------------------------------------------------------------ title page
\begin{center}
  {\LARGE\bfseries Sample Exam}\\[2mm]
  {\Large Medical Image Analysis}\\[6mm]
\end{center}

\vspace{4mm}
\noindent Student name:\hrulefill

\vspace{6mm}
\noindent Immatriculation number:\hrulefill

\vspace{10mm}

\begin{center}
  {\large\bfseries Exam Questions}
\end{center}

\noindent\textbf{Note:} All the questions have equal weight.

\vspace{4mm}

\begin{questions}

% =============================================================== Question 1 ===
\question
A measurement taken in an image corresponds to a real-world quantity only through the voxel size.
An anatomical structure is segmented in two scans of the same patient, acquired one year apart on
different scanners.

\begin{parts}

  \part In the \emph{baseline} scan the voxel size is $0.8 \times 0.8 \times 3.0$~mm$^3$ and the
  segmentation contains $4\,500$ voxels. What is the volume of the structure, in mm$^3$ and in
  cm$^3$?

  \begin{solution}
    Answer: $8\,640$~mm$^3$ $= 8.64$~cm$^3$.

    One voxel occupies $0.8 \times 0.8 \times 3.0 = 1.92$~mm$^3$, so
    $4\,500 \times 1.92 = 8\,640$~mm$^3$. There are $1\,000$~mm$^3$ in a cm$^3$.
  \end{solution}
\answerspace{2.4cm}

  \part In the \emph{follow-up} scan the voxel size is $0.5 \times 0.5 \times 2.0$~mm$^3$ and the
  segmentation contains $16\,000$ voxels. What is the volume now, and has the structure grown?

  \begin{solution}
    Answer: $8\,000$~mm$^3$ $= 8.00$~cm$^3$. The structure has \emph{not} grown --- it has shrunk
    slightly.

    A voxel now occupies $0.5 \times 0.5 \times 2.0 = 0.5$~mm$^3$, so
    $16\,000 \times 0.5 = 8\,000$~mm$^3$. Compared with the baseline this is a change of
    $(8\,000 - 8\,640)/8\,640 = -7.4\%$.
  \end{solution}
\answerspace{2.8cm}

  \part What would you have concluded if you had compared the raw voxel counts instead, and what
  does this tell you about comparing measurements across acquisitions?

  \begin{solution}
    Answer: you would have concluded that the structure grew by a factor of
    $16\,000 / 4\,500 \approx 3.6$, i.e.\ by more than $250\%$ --- the opposite of the truth, and
    wrong by a wide margin.

    A voxel count is not a measurement of anything physical. It becomes one only after being
    multiplied by the voxel volume, and the voxel volume differs between the two acquisitions here
    by almost a factor of four. Whenever measurements from different images are compared ---
    longitudinal studies, population statistics, normative distributions --- the real-world voxel
    size has to be taken into account, either by converting to physical units as here, or by
    resampling the images to a common voxel size first.
  \end{solution}
\answerspace{4.2cm}

\end{parts}

% =============================================================== Question 2 ===
\question
Deep learning libraries provide a layer called ``convolution'', but the operation they implement is
usually the closely related cross-correlation. For an image $I$ and a kernel $k$, which of the two
expressions below is the \emph{convolution} $I * k$? Please circle the correct choice.

\vspace{2mm}
\begin{center}
  $\displaystyle \sum_m \sum_n I(x-m,\, y-n)\, k(m,n)$
  \hspace{2.0cm}
  $\displaystyle \sum_m \sum_n I(x+m,\, y+n)\, k(m,n)$
\end{center}
\vspace{2mm}

\begin{solution}
  Answer: Left one.

  Convolution reverses one of the two functions before shifting it, which is what the minus signs
  encode; the right-hand expression has no flip and is the cross-correlation. The only thing the
  flip buys is commutativity, $I * k = k * I$, which is of little practical use in a neural network
  --- and since the kernel weights are learned anyway, a flipped kernel is just as learnable as an
  unflipped one. That is why most libraries implement cross-correlation and call it convolution,
  and why the distinction almost never matters in practice. It does matter when you take a kernel
  from a signal-processing textbook and expect it to behave identically.
\end{solution}
\answerspace{1.5cm}

% =============================================================== Question 3 ===
\question
A great many methods in this course have the same shape: a data term measuring agreement with the
observation, plus a regularisation term expressing which solutions are plausible,

\vspace{1mm}
\begin{center}
  $\displaystyle \min_\theta\ D(\text{data};\theta) \; + \; \lambda\, R(\theta).$
\end{center}
\vspace{1mm}

\noindent For each of the three methods below, please state what plays the role of the data term
and what plays the role of the regulariser.

\begin{parts}
  \part Total variation denoising.
  \begin{solution}
    Answer: data term $D = \|I - J\|_2^2$, the squared difference between the recovered image and
    the noisy observation. Regulariser $R = \|I\|_{TV} = \int \|\nabla I\|\,dx$, the total
    variation, which prefers images with minimal variation and therefore a piecewise-constant
    tendency.
  \end{solution}
\answerspace{2.6cm}

  \part Deformable intensity-based registration with linear-elastic regularisation.
  \begin{solution}
    Answer: data term $D = d(I,\, T_\theta \circ J)$, an intensity similarity metric such as the
    sum of squared differences between the target and the warped moving image. Regulariser
    $R = r(T_\theta^{-1})$, the elastic strain energy of the displacement field
    $u$, built from the Cauchy strain tensor $V(u) = (\nabla u + \nabla u^T)/2$, which prefers
    small, smooth deformations.
  \end{solution}
\answerspace{2.8cm}

  \part MAP segmentation with a Potts / Ising Markov random field prior.
  \begin{solution}
    Answer: data term is the unary term
    $\sum_x (I(x)-\mu_{c(x)})^T \Sigma^{-1}_{c(x)} (I(x)-\mu_{c(x)})$, i.e.\ the negative
    log-likelihood of the intensity at each pixel under the class it was assigned. Regulariser is
    the pairwise term $\lambda \sum_x \sum_{y \in G(x)} \delta(c(x) \neq c(y))$, which charges a
    fixed penalty for every neighbouring pair of pixels carrying different labels and therefore
    prefers spatially consistent labellings.
  \end{solution}
\answerspace{3.0cm}
\end{parts}

% =============================================================== Question 4 ===
\question
Please indicate whether the statement is true or false by circling the correct choice in each of
the following.

\begin{parts}
  \part \TFT\ \ Convolution is translation equivariant but not translation invariant.

  \part \TFF\ \ A ``same'' (padded) convolution with a $k \times k$ kernel and stride 1 reduces
  each spatial dimension by $k-1$.

  \part \TFF\ \ Fully connected layers are naturally translation invariant, which is why they are
  placed at the end of a classification network.

  \part \TFT\ \ For a recognition task two images differing only by a translation should give
  identical outputs, whereas for segmentation they should give the same output at a translated
  location.
\end{parts}

\begin{solution}
  (a) TRUE --- and the distinction is worth being precise about. Equivariance means
  $f(T(x)) = T(f(x))$: translate the input and the feature map translates with it. Invariance would
  mean $f(T(x)) = f(x)$, the output not changing at all, and convolution does \emph{not} give that.
  Pooling is what contributes a partial, local invariance.

  (b) FALSE --- that is the \emph{valid} convolution, $N_l = N_{l-1} - k_1 + 1$. The whole purpose
  of ``same'' convolution is to pad the boundary so the output has the same size as the input,
  $N_l = N_{l-1}$.

  (c) FALSE --- exactly backwards. Translation invariance is \emph{not} native to fully connected
  layers: translating the object produces completely different hidden activations, because every
  input position has its own separate weight. One can teach an FC network some invariance by
  augmenting the training data with random translations, but the lecture's verdict on that is
  ``not the most elegant way''.

  (d) TRUE --- and this is why the two properties are both wanted, in different places.
  Recognition wants invariance: a tumour is a tumour wherever it sits. Segmentation and detection
  want equivariance: move the structure and the output label map should move with it.
\end{solution}

% =============================================================== Question 5 ===
\question
A multi-head self-attention block has input dimension $D = 512$ and $H = 8$ heads, with per-head
dimensions $D_K = D_Q = D_V = 64$. Counting the projection matrices $W_{K_h}, W_{Q_h}, W_{V_h}$ of
every head together with the output projection $W_0$, and ignoring all biases, how many parameters
does the block have?

\begin{solution}
  Answer: $\mathbf{1\,048\,576}$ parameters.

  Per head, each of $W_{K_h}, W_{Q_h}, W_{V_h}$ has shape $D \times 64 = 512 \times 64 = 32\,768$
  entries, so a head contributes $3 \times 32\,768 = 98\,304$. Over 8 heads:
  $8 \times 98\,304 = 786\,432$.

  The heads are concatenated, giving $H \cdot D_V = 8 \times 64 = 512$ features, which the output
  projection maps back to $D$: $W_0 \in \mathbb{R}^{512 \times 512} = 262\,144$ entries.

  Total: $786\,432 + 262\,144 = 1\,048\,576$. Note that with $H \cdot D_V = D$, which is the usual
  choice, this is exactly $4D^2$ --- the cost of the block is quadratic in the model width and does
  not depend on the number of heads.
\end{solution}
\answerspace{4.0cm}

% =============================================================== Question 6 ===
\question
Once a Markov random field prior is added, segmentation can no longer be solved by deciding each
pixel independently. Please explain briefly why exact inference becomes hard, and describe one
approximate method presented in the lecture. You do not need to write down the equations.

\begin{solution}
  Answer: \emph{Why it is hard.} Without the MRF the posterior factorises over pixels, so each
  pixel's label can be chosen independently and the answer is found in one pass. The pairwise term
  couples neighbouring labels, so the posterior no longer factorises: the best label for a pixel
  depends on its neighbours' labels, which depend on theirs, and so on across the whole image.
  Computing the normalisation constant would require summing over every possible labelling of the
  image, of which there are exponentially many, and exact energy minimisation with categorical
  labels is NP-hard in 2D and above.

  \emph{An approximate method --- Gibbs sampling.} Sampling from the full posterior is infeasible,
  but sampling from the conditional distribution of a \emph{single} voxel given all the others is
  easy, because that conditional only involves the voxel's likelihood and its immediate
  neighbours, and normalising it means summing over the few possible labels of one voxel. So:
  initialise the labelling (for instance from the likelihood alone), then repeatedly visit each
  voxel and re-sample its label from that conditional. It is simple to implement and effective, but
  it is stochastic --- every run gives a different result --- the outcome depends on the
  initialisation and the parameters, and convergence can be slow.

  Iterated conditional modes and graph cuts are the other standard alternatives.
\end{solution}
\answerspace{7.5cm}

% =============================================================== Question 7 ===
\question
Which of the following statements on the differences between Active Shape Models (ASM) and Active
Appearance Models (AAM) is \textbf{inaccurate}?

\begin{choices}
  \choice ASM searches \emph{around} the current position, sampling along the surface normals of
  the model points, while AAM evaluates its fit \emph{at} the current position.
  \choice ASM minimises the distance between the model points and the corresponding points found
  in the image, whereas AAM minimises the difference between a synthesised model image and the
  target image.
  \choice AAM usually has a larger capture range than ASM, because the additional texture
  information gives it more to lock onto when the initialisation is poor.
  \choice AAM builds its texture model by warping every training image to the mean shape before
  sampling and normalising the intensities.
\end{choices}

\begin{solution}
  Answer: C.

  It is \textbf{ASM} that has the larger capture range, and (A) explains why: by sampling along the
  normals it actively looks a certain distance away from where the model currently sits, so it can
  find a boundary that is not underneath the model yet. AAM only compares its synthesised
  appearance with the image at the current position, so if the initialisation is too far off there
  is nothing to pull it in the right direction. ASM is also the faster of the two and locates
  feature points more accurately; what AAM buys in exchange is a better match to the image texture.
\end{solution}

% =============================================================== Question 8 ===
\question
Which of the following statements on non-linear transformation models is \textbf{incorrect}?

\begin{choices}
  \choice A naive parameterisation with one independent displacement vector per voxel can generate
  transformations that are physically implausible and even change the topology of the anatomy.
  \choice In a flow model each point follows the ordinary differential equation
  $dx/dt = v(x,t)$, and if the velocity field is smooth the resulting transformation is guaranteed
  to be a diffeomorphism.
  \choice A stationary velocity field has more parameters than a time-dependent one, which is why
  the scaling-and-squaring scheme is needed to make it tractable.
  \choice The diffusion model requires the displacement field to satisfy the Poisson equation,
  $\Delta u(x') + F(x') = 0$.
\end{choices}

\begin{solution}
  Answer: C.

  A stationary velocity field has \emph{fewer} parameters, not more: dropping the time dependence
  means storing one velocity vector per voxel instead of one per voxel \emph{per time point}. And
  scaling-and-squaring is not a workaround for a cost problem --- it is an efficient integration
  scheme that becomes available precisely \emph{because} the field is stationary. Both halves of
  the statement are inverted.

  The value of (B) is worth remembering separately: a diffeomorphism is smooth, invertible and has
  a smooth inverse, so a transformation built this way physically cannot tear or fold tissue. This
  is the guarantee that the naive parameterisation in (A) lacks, and it is the foundation of LDDMM.
\end{solution}

% =============================================================== Question 9 ===
\question
Which of the following statements on intensity normalisation is \textbf{incorrect}?

\begin{choices}
  \choice Min-max normalisation and matching the mean and standard deviation are both linear
  mappings of the intensities, and neither can change the contrast of an image.
  \choice Nyul's method matches the \emph{locations} of landmarks of the histograms, which aligns
  the distributions without making the two histograms identical.
  \choice Histogram matching exploits the spatial correlation between neighbouring pixels, which is
  what makes it robust to differences in the size of the imaged object.
  \choice Intensity variation is a particular problem for MRI because it has no absolute intensity
  scale, unlike CT where intensities are expressed in Hounsfield units.
\end{choices}

\begin{solution}
  Answer: C.

  Both halves are wrong. Histogram matching knows nothing at all about image content: it treats
  every pixel as an independent sample and uses \emph{no} spatial information or correlation
  between neighbours --- which is why matching image \emph{gradients} has been proposed as a
  remedy. And differences in object size are a genuine problem for it rather than something it is
  robust to: if one image contains much more of a given tissue than the other, a perfect histogram
  match is not even desirable, and mode matching is preferred.

  The other caveat worth knowing alongside this one is pathology: a large tumour changes the
  histogram, so matching a pathological image to a healthy one is hard, and robust statistics or
  outlier detection help.
\end{solution}

% ============================================================== Question 10 ===
\question
In patch-based label fusion the label at a target voxel is decided by
$L(x) = \arg\max_l \sum_i w_i\, \delta(\text{label}_i = l) \,/\, \sum_i w_i$, where the weights
$w_i$ come from the similarity between the patch around the target voxel and patches taken from
the atlas library. Which of the following statements about this scheme is correct?

\begin{choices}
  \choice The weights are computed from the spatial distance between the target voxel and the atlas
  voxel, not from the image content.
  \choice Only a coarse rigid or affine registration of the atlases is required, because
  correspondence is re-established locally by comparing patches inside a search window around each
  voxel.
  \choice The scheme requires exactly one atlas; with several atlases the weights no longer sum to
  one and the decision becomes ill-defined.
  \choice Because the weighting is borrowed from non-local means, the scheme can only be applied to
  denoising and not to the propagation of labels.
  \choice The search window has to cover the whole image, since restricting it would bias the
  weights.
\end{choices}

\begin{solution}
  Answer: B.

  This is the point of patch-based methods. In classical multi-atlas segmentation the
  correspondence between the atlas and the patient is whatever the registration says it is, so the
  result is only as good as a non-linear registration --- which is expensive and can fail. Patch
  matching replaces that: after a coarse alignment, each target voxel looks around a local search
  window for atlas patches that actually resemble its own neighbourhood, and takes their labels
  weighted by that resemblance. Correspondence is decided by image content at the point of use
  rather than by a global transformation computed in advance.

  (E) is the reverse of the truth: restricting the search window is standard practice, for cost, and
  it also encodes the sensible prior that the correct correspondence is somewhere nearby.
\end{solution}

% ============================================================== Question 11 ===
\question
A U-Net for segmentation takes $128 \times 128$ inputs with one channel. All convolutions are
$3 \times 3$ with ``same'' padding; each encoder level ends with $2 \times 2$ max pooling of stride
2, and each decoder level begins by upsampling by a factor of 2. The encoder levels produce 32, 64
and 128 channels respectively, and the bottleneck produces 256 channels. Each skip connection
concatenates the encoder features onto the upsampled decoder features at the matching level. Circle
True or False:

\begin{parts}
  \part \TFT\ \ The feature maps at the bottleneck have a spatial size of $16 \times 16$.

  \part \TFF\ \ Because all convolutions use ``same'' padding, no spatial information is lost in
  the encoder, and the skip connections are therefore redundant.

  \part \TFT\ \ After upsampling the bottleneck to $32 \times 32$ and concatenating the skip
  connection from the matching encoder level, the first decoder convolution at that level receives
  $256 + 128 = 384$ input channels.

  \part \TFF\ \ A $3 \times 3$ convolution mapping those 384 input channels to 128 output channels
  has $3 \cdot 3 \cdot 384 = 3\,456$ weights, excluding biases.
\end{parts}

\begin{solution}
  (a) TRUE --- three poolings, each halving: $128 \rightarrow 64 \rightarrow 32 \rightarrow 16$.
  The ``same'' convolutions do not change the size, so only the pooling does.

  (b) FALSE --- and this is the point of the architecture. It is the \emph{pooling}, not the
  convolutions, that destroys spatial detail: by the bottleneck the feature maps are $16 \times 16$
  and the fine boundary information is gone. The skip connections exist to carry the
  high-resolution features from the encoder across to the decoder, so the network can combine the
  large receptive field of the deep layers with the spatial precision of the shallow ones.

  (c) TRUE --- concatenation adds the channel counts: 256 upsampled from the bottleneck plus the
  128 arriving through the skip connection.

  (d) FALSE --- the count leaves out the output channels. A convolution kernel connecting every
  input channel to every output channel has $k_1 \cdot k_2 \cdot C_{in} \cdot C_{out}$ weights, so
  the correct number is $3 \cdot 3 \cdot 384 \cdot 128 = 442\,368$ --- 128 times the stated figure.
  Note also that this is where much of a U-Net's parameter budget goes: the concatenation makes the
  decoder convolutions wide.
\end{solution}

% ============================================================== Question 12 ===
\question
Which of the following statements on interpretability and its uses are correct (multiple answers
are possible)?

\begin{choices}
  \choice Saliency maps change over the course of training, and can be used as a kind of
  fingerprint of how a model responds to its inputs.
  \choice In interpretability-guided active learning, a sample is treated as informative when the
  model appears to struggle with it; easy samples tend to produce clean, focused saliency.
  \choice A model-agnostic method such as LIME explains a prediction by fitting a simpler,
  interpretable model that approximates the classifier in a neighbourhood of the input, weighted by
  a proximity measure.
  \choice Class-prototype maximisation returns the image from the training set that most strongly
  activates a given class.
  \choice Grad-CAM can be applied without any access to the internals of the model, since it only
  requires the class scores.
\end{choices}

\begin{solution}
  Answer: (A), (B) and (C) are correct.

  (D) misdescribes the method. Class-prototype maximisation performs gradient \emph{ascent on the
  input} to synthesise a pattern that maximises a class activation; the result is not a training
  image and indeed is not necessarily a realistic image at all. Regularising it with a generative
  prior is what makes the synthesised inputs look natural.

  (E) puts Grad-CAM on the wrong side of the taxonomy. It is a \emph{gradient}-based method, and it
  needs the gradients of the class score with respect to the convolutional feature maps --- so it
  requires access to the internals and is a white-box method. LIME is the model-agnostic,
  black-box one.
\end{solution}

\end{questions}

\end{document}
"
%% Compile the blank version (no answers):  pdflatex Sample_Exam_5.tex
%%
\ifx\withanswers\undefined
  \documentclass[11pt]{exam}
\else
  \documentclass[11pt,answers]{exam}
\fi

\usepackage[margin=2.4cm]{geometry}
\usepackage{amsmath,amssymb}
\usepackage{array}
\usepackage[T1]{fontenc}

% ---------------------------------------------------------------- page style
\pagestyle{foot}
\firstpageheader{}{}{}
\firstpagefooter{}{}{}
\runningheader{}{}{}
\runningfooter{}{}{Page \thepage}

% -------------------------------------------------------------- multiplechoice
% Reproduce the "(A) (B) (C) ..." labelling of the original exam.
\renewcommand{\choicelabel}{$\bigcirc$~(\Alph{choice})}

% ---------------------------------------------------------- true/false macros
% \TFT -> the statement is TRUE, \TFF -> the statement is FALSE.
% In the answer version the correct word is set in capitals and bold.
\newcommand{\TFT}{\ifprintanswers\textbf{TRUE}\hspace{1.1em}False\else True\hspace{1.1em}False\fi}
\newcommand{\TFF}{True\hspace{1.1em}\ifprintanswers\textbf{FALSE}\else False\fi}

% ------------------------------------------------------------- writing space
% Reserves room to write in the blank version; collapses in the answer version.
\newcommand{\answerspace}[1]{\ifprintanswers\else\vspace{#1}\fi}

\begin{document}

% ------------------------------------------------------------------ title page
\begin{center}
  {\LARGE\bfseries Sample Exam}\\[2mm]
  {\Large Medical Image Analysis}\\[6mm]
\end{center}

\vspace{4mm}
\noindent Student name:\hrulefill

\vspace{6mm}
\noindent Immatriculation number:\hrulefill

\vspace{10mm}

\begin{center}
  {\large\bfseries Exam Questions}
\end{center}

\noindent\textbf{Note:} All the questions have equal weight.

\vspace{4mm}

\begin{questions}

% =============================================================== Question 1 ===
\question
A measurement taken in an image corresponds to a real-world quantity only through the voxel size.
An anatomical structure is segmented in two scans of the same patient, acquired one year apart on
different scanners.

\begin{parts}

  \part In the \emph{baseline} scan the voxel size is $0.8 \times 0.8 \times 3.0$~mm$^3$ and the
  segmentation contains $4\,500$ voxels. What is the volume of the structure, in mm$^3$ and in
  cm$^3$?

  \begin{solution}
    Answer: $8\,640$~mm$^3$ $= 8.64$~cm$^3$.

    One voxel occupies $0.8 \times 0.8 \times 3.0 = 1.92$~mm$^3$, so
    $4\,500 \times 1.92 = 8\,640$~mm$^3$. There are $1\,000$~mm$^3$ in a cm$^3$.
  \end{solution}
\answerspace{2.4cm}

  \part In the \emph{follow-up} scan the voxel size is $0.5 \times 0.5 \times 2.0$~mm$^3$ and the
  segmentation contains $16\,000$ voxels. What is the volume now, and has the structure grown?

  \begin{solution}
    Answer: $8\,000$~mm$^3$ $= 8.00$~cm$^3$. The structure has \emph{not} grown --- it has shrunk
    slightly.

    A voxel now occupies $0.5 \times 0.5 \times 2.0 = 0.5$~mm$^3$, so
    $16\,000 \times 0.5 = 8\,000$~mm$^3$. Compared with the baseline this is a change of
    $(8\,000 - 8\,640)/8\,640 = -7.4\%$.
  \end{solution}
\answerspace{2.8cm}

  \part What would you have concluded if you had compared the raw voxel counts instead, and what
  does this tell you about comparing measurements across acquisitions?

  \begin{solution}
    Answer: you would have concluded that the structure grew by a factor of
    $16\,000 / 4\,500 \approx 3.6$, i.e.\ by more than $250\%$ --- the opposite of the truth, and
    wrong by a wide margin.

    A voxel count is not a measurement of anything physical. It becomes one only after being
    multiplied by the voxel volume, and the voxel volume differs between the two acquisitions here
    by almost a factor of four. Whenever measurements from different images are compared ---
    longitudinal studies, population statistics, normative distributions --- the real-world voxel
    size has to be taken into account, either by converting to physical units as here, or by
    resampling the images to a common voxel size first.
  \end{solution}
\answerspace{4.2cm}

\end{parts}

% =============================================================== Question 2 ===
\question
Deep learning libraries provide a layer called ``convolution'', but the operation they implement is
usually the closely related cross-correlation. For an image $I$ and a kernel $k$, which of the two
expressions below is the \emph{convolution} $I * k$? Please circle the correct choice.

\vspace{2mm}
\begin{center}
  $\displaystyle \sum_m \sum_n I(x-m,\, y-n)\, k(m,n)$
  \hspace{2.0cm}
  $\displaystyle \sum_m \sum_n I(x+m,\, y+n)\, k(m,n)$
\end{center}
\vspace{2mm}

\begin{solution}
  Answer: Left one.

  Convolution reverses one of the two functions before shifting it, which is what the minus signs
  encode; the right-hand expression has no flip and is the cross-correlation. The only thing the
  flip buys is commutativity, $I * k = k * I$, which is of little practical use in a neural network
  --- and since the kernel weights are learned anyway, a flipped kernel is just as learnable as an
  unflipped one. That is why most libraries implement cross-correlation and call it convolution,
  and why the distinction almost never matters in practice. It does matter when you take a kernel
  from a signal-processing textbook and expect it to behave identically.
\end{solution}
\answerspace{1.5cm}

% =============================================================== Question 3 ===
\question
A great many methods in this course have the same shape: a data term measuring agreement with the
observation, plus a regularisation term expressing which solutions are plausible,

\vspace{1mm}
\begin{center}
  $\displaystyle \min_\theta\ D(\text{data};\theta) \; + \; \lambda\, R(\theta).$
\end{center}
\vspace{1mm}

\noindent For each of the three methods below, please state what plays the role of the data term
and what plays the role of the regulariser.

\begin{parts}
  \part Total variation denoising.
  \begin{solution}
    Answer: data term $D = \|I - J\|_2^2$, the squared difference between the recovered image and
    the noisy observation. Regulariser $R = \|I\|_{TV} = \int \|\nabla I\|\,dx$, the total
    variation, which prefers images with minimal variation and therefore a piecewise-constant
    tendency.
  \end{solution}
\answerspace{2.6cm}

  \part Deformable intensity-based registration with linear-elastic regularisation.
  \begin{solution}
    Answer: data term $D = d(I,\, T_\theta \circ J)$, an intensity similarity metric such as the
    sum of squared differences between the target and the warped moving image. Regulariser
    $R = r(T_\theta^{-1})$, the elastic strain energy of the displacement field
    $u$, built from the Cauchy strain tensor $V(u) = (\nabla u + \nabla u^T)/2$, which prefers
    small, smooth deformations.
  \end{solution}
\answerspace{2.8cm}

  \part MAP segmentation with a Potts / Ising Markov random field prior.
  \begin{solution}
    Answer: data term is the unary term
    $\sum_x (I(x)-\mu_{c(x)})^T \Sigma^{-1}_{c(x)} (I(x)-\mu_{c(x)})$, i.e.\ the negative
    log-likelihood of the intensity at each pixel under the class it was assigned. Regulariser is
    the pairwise term $\lambda \sum_x \sum_{y \in G(x)} \delta(c(x) \neq c(y))$, which charges a
    fixed penalty for every neighbouring pair of pixels carrying different labels and therefore
    prefers spatially consistent labellings.
  \end{solution}
\answerspace{3.0cm}
\end{parts}

% =============================================================== Question 4 ===
\question
Please indicate whether the statement is true or false by circling the correct choice in each of
the following.

\begin{parts}
  \part \TFT\ \ Convolution is translation equivariant but not translation invariant.

  \part \TFF\ \ A ``same'' (padded) convolution with a $k \times k$ kernel and stride 1 reduces
  each spatial dimension by $k-1$.

  \part \TFF\ \ Fully connected layers are naturally translation invariant, which is why they are
  placed at the end of a classification network.

  \part \TFT\ \ For a recognition task two images differing only by a translation should give
  identical outputs, whereas for segmentation they should give the same output at a translated
  location.
\end{parts}

\begin{solution}
  (a) TRUE --- and the distinction is worth being precise about. Equivariance means
  $f(T(x)) = T(f(x))$: translate the input and the feature map translates with it. Invariance would
  mean $f(T(x)) = f(x)$, the output not changing at all, and convolution does \emph{not} give that.
  Pooling is what contributes a partial, local invariance.

  (b) FALSE --- that is the \emph{valid} convolution, $N_l = N_{l-1} - k_1 + 1$. The whole purpose
  of ``same'' convolution is to pad the boundary so the output has the same size as the input,
  $N_l = N_{l-1}$.

  (c) FALSE --- exactly backwards. Translation invariance is \emph{not} native to fully connected
  layers: translating the object produces completely different hidden activations, because every
  input position has its own separate weight. One can teach an FC network some invariance by
  augmenting the training data with random translations, but the lecture's verdict on that is
  ``not the most elegant way''.

  (d) TRUE --- and this is why the two properties are both wanted, in different places.
  Recognition wants invariance: a tumour is a tumour wherever it sits. Segmentation and detection
  want equivariance: move the structure and the output label map should move with it.
\end{solution}

% =============================================================== Question 5 ===
\question
A multi-head self-attention block has input dimension $D = 512$ and $H = 8$ heads, with per-head
dimensions $D_K = D_Q = D_V = 64$. Counting the projection matrices $W_{K_h}, W_{Q_h}, W_{V_h}$ of
every head together with the output projection $W_0$, and ignoring all biases, how many parameters
does the block have?

\begin{solution}
  Answer: $\mathbf{1\,048\,576}$ parameters.

  Per head, each of $W_{K_h}, W_{Q_h}, W_{V_h}$ has shape $D \times 64 = 512 \times 64 = 32\,768$
  entries, so a head contributes $3 \times 32\,768 = 98\,304$. Over 8 heads:
  $8 \times 98\,304 = 786\,432$.

  The heads are concatenated, giving $H \cdot D_V = 8 \times 64 = 512$ features, which the output
  projection maps back to $D$: $W_0 \in \mathbb{R}^{512 \times 512} = 262\,144$ entries.

  Total: $786\,432 + 262\,144 = 1\,048\,576$. Note that with $H \cdot D_V = D$, which is the usual
  choice, this is exactly $4D^2$ --- the cost of the block is quadratic in the model width and does
  not depend on the number of heads.
\end{solution}
\answerspace{4.0cm}

% =============================================================== Question 6 ===
\question
Once a Markov random field prior is added, segmentation can no longer be solved by deciding each
pixel independently. Please explain briefly why exact inference becomes hard, and describe one
approximate method presented in the lecture. You do not need to write down the equations.

\begin{solution}
  Answer: \emph{Why it is hard.} Without the MRF the posterior factorises over pixels, so each
  pixel's label can be chosen independently and the answer is found in one pass. The pairwise term
  couples neighbouring labels, so the posterior no longer factorises: the best label for a pixel
  depends on its neighbours' labels, which depend on theirs, and so on across the whole image.
  Computing the normalisation constant would require summing over every possible labelling of the
  image, of which there are exponentially many, and exact energy minimisation with categorical
  labels is NP-hard in 2D and above.

  \emph{An approximate method --- Gibbs sampling.} Sampling from the full posterior is infeasible,
  but sampling from the conditional distribution of a \emph{single} voxel given all the others is
  easy, because that conditional only involves the voxel's likelihood and its immediate
  neighbours, and normalising it means summing over the few possible labels of one voxel. So:
  initialise the labelling (for instance from the likelihood alone), then repeatedly visit each
  voxel and re-sample its label from that conditional. It is simple to implement and effective, but
  it is stochastic --- every run gives a different result --- the outcome depends on the
  initialisation and the parameters, and convergence can be slow.

  Iterated conditional modes and graph cuts are the other standard alternatives.
\end{solution}
\answerspace{7.5cm}

% =============================================================== Question 7 ===
\question
Which of the following statements on the differences between Active Shape Models (ASM) and Active
Appearance Models (AAM) is \textbf{inaccurate}?

\begin{choices}
  \choice ASM searches \emph{around} the current position, sampling along the surface normals of
  the model points, while AAM evaluates its fit \emph{at} the current position.
  \choice ASM minimises the distance between the model points and the corresponding points found
  in the image, whereas AAM minimises the difference between a synthesised model image and the
  target image.
  \choice AAM usually has a larger capture range than ASM, because the additional texture
  information gives it more to lock onto when the initialisation is poor.
  \choice AAM builds its texture model by warping every training image to the mean shape before
  sampling and normalising the intensities.
\end{choices}

\begin{solution}
  Answer: C.

  It is \textbf{ASM} that has the larger capture range, and (A) explains why: by sampling along the
  normals it actively looks a certain distance away from where the model currently sits, so it can
  find a boundary that is not underneath the model yet. AAM only compares its synthesised
  appearance with the image at the current position, so if the initialisation is too far off there
  is nothing to pull it in the right direction. ASM is also the faster of the two and locates
  feature points more accurately; what AAM buys in exchange is a better match to the image texture.
\end{solution}

% =============================================================== Question 8 ===
\question
Which of the following statements on non-linear transformation models is \textbf{incorrect}?

\begin{choices}
  \choice A naive parameterisation with one independent displacement vector per voxel can generate
  transformations that are physically implausible and even change the topology of the anatomy.
  \choice In a flow model each point follows the ordinary differential equation
  $dx/dt = v(x,t)$, and if the velocity field is smooth the resulting transformation is guaranteed
  to be a diffeomorphism.
  \choice A stationary velocity field has more parameters than a time-dependent one, which is why
  the scaling-and-squaring scheme is needed to make it tractable.
  \choice The diffusion model requires the displacement field to satisfy the Poisson equation,
  $\Delta u(x') + F(x') = 0$.
\end{choices}

\begin{solution}
  Answer: C.

  A stationary velocity field has \emph{fewer} parameters, not more: dropping the time dependence
  means storing one velocity vector per voxel instead of one per voxel \emph{per time point}. And
  scaling-and-squaring is not a workaround for a cost problem --- it is an efficient integration
  scheme that becomes available precisely \emph{because} the field is stationary. Both halves of
  the statement are inverted.

  The value of (B) is worth remembering separately: a diffeomorphism is smooth, invertible and has
  a smooth inverse, so a transformation built this way physically cannot tear or fold tissue. This
  is the guarantee that the naive parameterisation in (A) lacks, and it is the foundation of LDDMM.
\end{solution}

% =============================================================== Question 9 ===
\question
Which of the following statements on intensity normalisation is \textbf{incorrect}?

\begin{choices}
  \choice Min-max normalisation and matching the mean and standard deviation are both linear
  mappings of the intensities, and neither can change the contrast of an image.
  \choice Nyul's method matches the \emph{locations} of landmarks of the histograms, which aligns
  the distributions without making the two histograms identical.
  \choice Histogram matching exploits the spatial correlation between neighbouring pixels, which is
  what makes it robust to differences in the size of the imaged object.
  \choice Intensity variation is a particular problem for MRI because it has no absolute intensity
  scale, unlike CT where intensities are expressed in Hounsfield units.
\end{choices}

\begin{solution}
  Answer: C.

  Both halves are wrong. Histogram matching knows nothing at all about image content: it treats
  every pixel as an independent sample and uses \emph{no} spatial information or correlation
  between neighbours --- which is why matching image \emph{gradients} has been proposed as a
  remedy. And differences in object size are a genuine problem for it rather than something it is
  robust to: if one image contains much more of a given tissue than the other, a perfect histogram
  match is not even desirable, and mode matching is preferred.

  The other caveat worth knowing alongside this one is pathology: a large tumour changes the
  histogram, so matching a pathological image to a healthy one is hard, and robust statistics or
  outlier detection help.
\end{solution}

% ============================================================== Question 10 ===
\question
In patch-based label fusion the label at a target voxel is decided by
$L(x) = \arg\max_l \sum_i w_i\, \delta(\text{label}_i = l) \,/\, \sum_i w_i$, where the weights
$w_i$ come from the similarity between the patch around the target voxel and patches taken from
the atlas library. Which of the following statements about this scheme is correct?

\begin{choices}
  \choice The weights are computed from the spatial distance between the target voxel and the atlas
  voxel, not from the image content.
  \choice Only a coarse rigid or affine registration of the atlases is required, because
  correspondence is re-established locally by comparing patches inside a search window around each
  voxel.
  \choice The scheme requires exactly one atlas; with several atlases the weights no longer sum to
  one and the decision becomes ill-defined.
  \choice Because the weighting is borrowed from non-local means, the scheme can only be applied to
  denoising and not to the propagation of labels.
  \choice The search window has to cover the whole image, since restricting it would bias the
  weights.
\end{choices}

\begin{solution}
  Answer: B.

  This is the point of patch-based methods. In classical multi-atlas segmentation the
  correspondence between the atlas and the patient is whatever the registration says it is, so the
  result is only as good as a non-linear registration --- which is expensive and can fail. Patch
  matching replaces that: after a coarse alignment, each target voxel looks around a local search
  window for atlas patches that actually resemble its own neighbourhood, and takes their labels
  weighted by that resemblance. Correspondence is decided by image content at the point of use
  rather than by a global transformation computed in advance.

  (E) is the reverse of the truth: restricting the search window is standard practice, for cost, and
  it also encodes the sensible prior that the correct correspondence is somewhere nearby.
\end{solution}

% ============================================================== Question 11 ===
\question
A U-Net for segmentation takes $128 \times 128$ inputs with one channel. All convolutions are
$3 \times 3$ with ``same'' padding; each encoder level ends with $2 \times 2$ max pooling of stride
2, and each decoder level begins by upsampling by a factor of 2. The encoder levels produce 32, 64
and 128 channels respectively, and the bottleneck produces 256 channels. Each skip connection
concatenates the encoder features onto the upsampled decoder features at the matching level. Circle
True or False:

\begin{parts}
  \part \TFT\ \ The feature maps at the bottleneck have a spatial size of $16 \times 16$.

  \part \TFF\ \ Because all convolutions use ``same'' padding, no spatial information is lost in
  the encoder, and the skip connections are therefore redundant.

  \part \TFT\ \ After upsampling the bottleneck to $32 \times 32$ and concatenating the skip
  connection from the matching encoder level, the first decoder convolution at that level receives
  $256 + 128 = 384$ input channels.

  \part \TFF\ \ A $3 \times 3$ convolution mapping those 384 input channels to 128 output channels
  has $3 \cdot 3 \cdot 384 = 3\,456$ weights, excluding biases.
\end{parts}

\begin{solution}
  (a) TRUE --- three poolings, each halving: $128 \rightarrow 64 \rightarrow 32 \rightarrow 16$.
  The ``same'' convolutions do not change the size, so only the pooling does.

  (b) FALSE --- and this is the point of the architecture. It is the \emph{pooling}, not the
  convolutions, that destroys spatial detail: by the bottleneck the feature maps are $16 \times 16$
  and the fine boundary information is gone. The skip connections exist to carry the
  high-resolution features from the encoder across to the decoder, so the network can combine the
  large receptive field of the deep layers with the spatial precision of the shallow ones.

  (c) TRUE --- concatenation adds the channel counts: 256 upsampled from the bottleneck plus the
  128 arriving through the skip connection.

  (d) FALSE --- the count leaves out the output channels. A convolution kernel connecting every
  input channel to every output channel has $k_1 \cdot k_2 \cdot C_{in} \cdot C_{out}$ weights, so
  the correct number is $3 \cdot 3 \cdot 384 \cdot 128 = 442\,368$ --- 128 times the stated figure.
  Note also that this is where much of a U-Net's parameter budget goes: the concatenation makes the
  decoder convolutions wide.
\end{solution}

% ============================================================== Question 12 ===
\question
Which of the following statements on interpretability and its uses are correct (multiple answers
are possible)?

\begin{choices}
  \choice Saliency maps change over the course of training, and can be used as a kind of
  fingerprint of how a model responds to its inputs.
  \choice In interpretability-guided active learning, a sample is treated as informative when the
  model appears to struggle with it; easy samples tend to produce clean, focused saliency.
  \choice A model-agnostic method such as LIME explains a prediction by fitting a simpler,
  interpretable model that approximates the classifier in a neighbourhood of the input, weighted by
  a proximity measure.
  \choice Class-prototype maximisation returns the image from the training set that most strongly
  activates a given class.
  \choice Grad-CAM can be applied without any access to the internals of the model, since it only
  requires the class scores.
\end{choices}

\begin{solution}
  Answer: (A), (B) and (C) are correct.

  (D) misdescribes the method. Class-prototype maximisation performs gradient \emph{ascent on the
  input} to synthesise a pattern that maximises a class activation; the result is not a training
  image and indeed is not necessarily a realistic image at all. Regularising it with a generative
  prior is what makes the synthesised inputs look natural.

  (E) puts Grad-CAM on the wrong side of the taxonomy. It is a \emph{gradient}-based method, and it
  needs the gradients of the class score with respect to the convolutional feature maps --- so it
  requires access to the internals and is a white-box method. LIME is the model-agnostic,
  black-box one.
\end{solution}

\end{questions}

\end{document}
"
%% Compile the blank version (no answers):  pdflatex Sample_Exam_5.tex
%%
\ifx\withanswers\undefined
  \documentclass[11pt]{exam}
\else
  \documentclass[11pt,answers]{exam}
\fi

\usepackage[margin=2.4cm]{geometry}
\usepackage{amsmath,amssymb}
\usepackage{array}
\usepackage[T1]{fontenc}

% ---------------------------------------------------------------- page style
\pagestyle{foot}
\firstpageheader{}{}{}
\firstpagefooter{}{}{}
\runningheader{}{}{}
\runningfooter{}{}{Page \thepage}

% -------------------------------------------------------------- multiplechoice
% Reproduce the "(A) (B) (C) ..." labelling of the original exam.
\renewcommand{\choicelabel}{$\bigcirc$~(\Alph{choice})}

% ---------------------------------------------------------- true/false macros
% \TFT -> the statement is TRUE, \TFF -> the statement is FALSE.
% In the answer version the correct word is set in capitals and bold.
\newcommand{\TFT}{\ifprintanswers\textbf{TRUE}\hspace{1.1em}False\else True\hspace{1.1em}False\fi}
\newcommand{\TFF}{True\hspace{1.1em}\ifprintanswers\textbf{FALSE}\else False\fi}

% ------------------------------------------------------------- writing space
% Reserves room to write in the blank version; collapses in the answer version.
\newcommand{\answerspace}[1]{\ifprintanswers\else\vspace{#1}\fi}

\begin{document}

% ------------------------------------------------------------------ title page
\begin{center}
  {\LARGE\bfseries Sample Exam}\\[2mm]
  {\Large Medical Image Analysis}\\[6mm]
\end{center}

\vspace{4mm}
\noindent Student name:\hrulefill

\vspace{6mm}
\noindent Immatriculation number:\hrulefill

\vspace{10mm}

\begin{center}
  {\large\bfseries Exam Questions}
\end{center}

\noindent\textbf{Note:} All the questions have equal weight.

\vspace{4mm}

\begin{questions}

% =============================================================== Question 1 ===
\question
A measurement taken in an image corresponds to a real-world quantity only through the voxel size.
An anatomical structure is segmented in two scans of the same patient, acquired one year apart on
different scanners.

\begin{parts}

  \part In the \emph{baseline} scan the voxel size is $0.8 \times 0.8 \times 3.0$~mm$^3$ and the
  segmentation contains $4\,500$ voxels. What is the volume of the structure, in mm$^3$ and in
  cm$^3$?

  \begin{solution}
    Answer: $8\,640$~mm$^3$ $= 8.64$~cm$^3$.

    One voxel occupies $0.8 \times 0.8 \times 3.0 = 1.92$~mm$^3$, so
    $4\,500 \times 1.92 = 8\,640$~mm$^3$. There are $1\,000$~mm$^3$ in a cm$^3$.
  \end{solution}
\answerspace{2.4cm}

  \part In the \emph{follow-up} scan the voxel size is $0.5 \times 0.5 \times 2.0$~mm$^3$ and the
  segmentation contains $16\,000$ voxels. What is the volume now, and has the structure grown?

  \begin{solution}
    Answer: $8\,000$~mm$^3$ $= 8.00$~cm$^3$. The structure has \emph{not} grown --- it has shrunk
    slightly.

    A voxel now occupies $0.5 \times 0.5 \times 2.0 = 0.5$~mm$^3$, so
    $16\,000 \times 0.5 = 8\,000$~mm$^3$. Compared with the baseline this is a change of
    $(8\,000 - 8\,640)/8\,640 = -7.4\%$.
  \end{solution}
\answerspace{2.8cm}

  \part What would you have concluded if you had compared the raw voxel counts instead, and what
  does this tell you about comparing measurements across acquisitions?

  \begin{solution}
    Answer: you would have concluded that the structure grew by a factor of
    $16\,000 / 4\,500 \approx 3.6$, i.e.\ by more than $250\%$ --- the opposite of the truth, and
    wrong by a wide margin.

    A voxel count is not a measurement of anything physical. It becomes one only after being
    multiplied by the voxel volume, and the voxel volume differs between the two acquisitions here
    by almost a factor of four. Whenever measurements from different images are compared ---
    longitudinal studies, population statistics, normative distributions --- the real-world voxel
    size has to be taken into account, either by converting to physical units as here, or by
    resampling the images to a common voxel size first.
  \end{solution}
\answerspace{4.2cm}

\end{parts}

% =============================================================== Question 2 ===
\question
Deep learning libraries provide a layer called ``convolution'', but the operation they implement is
usually the closely related cross-correlation. For an image $I$ and a kernel $k$, which of the two
expressions below is the \emph{convolution} $I * k$? Please circle the correct choice.

\vspace{2mm}
\begin{center}
  $\displaystyle \sum_m \sum_n I(x-m,\, y-n)\, k(m,n)$
  \hspace{2.0cm}
  $\displaystyle \sum_m \sum_n I(x+m,\, y+n)\, k(m,n)$
\end{center}
\vspace{2mm}

\begin{solution}
  Answer: Left one.

  Convolution reverses one of the two functions before shifting it, which is what the minus signs
  encode; the right-hand expression has no flip and is the cross-correlation. The only thing the
  flip buys is commutativity, $I * k = k * I$, which is of little practical use in a neural network
  --- and since the kernel weights are learned anyway, a flipped kernel is just as learnable as an
  unflipped one. That is why most libraries implement cross-correlation and call it convolution,
  and why the distinction almost never matters in practice. It does matter when you take a kernel
  from a signal-processing textbook and expect it to behave identically.
\end{solution}
\answerspace{1.5cm}

% =============================================================== Question 3 ===
\question
A great many methods in this course have the same shape: a data term measuring agreement with the
observation, plus a regularisation term expressing which solutions are plausible,

\vspace{1mm}
\begin{center}
  $\displaystyle \min_\theta\ D(\text{data};\theta) \; + \; \lambda\, R(\theta).$
\end{center}
\vspace{1mm}

\noindent For each of the three methods below, please state what plays the role of the data term
and what plays the role of the regulariser.

\begin{parts}
  \part Total variation denoising.
  \begin{solution}
    Answer: data term $D = \|I - J\|_2^2$, the squared difference between the recovered image and
    the noisy observation. Regulariser $R = \|I\|_{TV} = \int \|\nabla I\|\,dx$, the total
    variation, which prefers images with minimal variation and therefore a piecewise-constant
    tendency.
  \end{solution}
\answerspace{2.6cm}

  \part Deformable intensity-based registration with linear-elastic regularisation.
  \begin{solution}
    Answer: data term $D = d(I,\, T_\theta \circ J)$, an intensity similarity metric such as the
    sum of squared differences between the target and the warped moving image. Regulariser
    $R = r(T_\theta^{-1})$, the elastic strain energy of the displacement field
    $u$, built from the Cauchy strain tensor $V(u) = (\nabla u + \nabla u^T)/2$, which prefers
    small, smooth deformations.
  \end{solution}
\answerspace{2.8cm}

  \part MAP segmentation with a Potts / Ising Markov random field prior.
  \begin{solution}
    Answer: data term is the unary term
    $\sum_x (I(x)-\mu_{c(x)})^T \Sigma^{-1}_{c(x)} (I(x)-\mu_{c(x)})$, i.e.\ the negative
    log-likelihood of the intensity at each pixel under the class it was assigned. Regulariser is
    the pairwise term $\lambda \sum_x \sum_{y \in G(x)} \delta(c(x) \neq c(y))$, which charges a
    fixed penalty for every neighbouring pair of pixels carrying different labels and therefore
    prefers spatially consistent labellings.
  \end{solution}
\answerspace{3.0cm}
\end{parts}

% =============================================================== Question 4 ===
\question
Please indicate whether the statement is true or false by circling the correct choice in each of
the following.

\begin{parts}
  \part \TFT\ \ Convolution is translation equivariant but not translation invariant.

  \part \TFF\ \ A ``same'' (padded) convolution with a $k \times k$ kernel and stride 1 reduces
  each spatial dimension by $k-1$.

  \part \TFF\ \ Fully connected layers are naturally translation invariant, which is why they are
  placed at the end of a classification network.

  \part \TFT\ \ For a recognition task two images differing only by a translation should give
  identical outputs, whereas for segmentation they should give the same output at a translated
  location.
\end{parts}

\begin{solution}
  (a) TRUE --- and the distinction is worth being precise about. Equivariance means
  $f(T(x)) = T(f(x))$: translate the input and the feature map translates with it. Invariance would
  mean $f(T(x)) = f(x)$, the output not changing at all, and convolution does \emph{not} give that.
  Pooling is what contributes a partial, local invariance.

  (b) FALSE --- that is the \emph{valid} convolution, $N_l = N_{l-1} - k_1 + 1$. The whole purpose
  of ``same'' convolution is to pad the boundary so the output has the same size as the input,
  $N_l = N_{l-1}$.

  (c) FALSE --- exactly backwards. Translation invariance is \emph{not} native to fully connected
  layers: translating the object produces completely different hidden activations, because every
  input position has its own separate weight. One can teach an FC network some invariance by
  augmenting the training data with random translations, but the lecture's verdict on that is
  ``not the most elegant way''.

  (d) TRUE --- and this is why the two properties are both wanted, in different places.
  Recognition wants invariance: a tumour is a tumour wherever it sits. Segmentation and detection
  want equivariance: move the structure and the output label map should move with it.
\end{solution}

% =============================================================== Question 5 ===
\question
A multi-head self-attention block has input dimension $D = 512$ and $H = 8$ heads, with per-head
dimensions $D_K = D_Q = D_V = 64$. Counting the projection matrices $W_{K_h}, W_{Q_h}, W_{V_h}$ of
every head together with the output projection $W_0$, and ignoring all biases, how many parameters
does the block have?

\begin{solution}
  Answer: $\mathbf{1\,048\,576}$ parameters.

  Per head, each of $W_{K_h}, W_{Q_h}, W_{V_h}$ has shape $D \times 64 = 512 \times 64 = 32\,768$
  entries, so a head contributes $3 \times 32\,768 = 98\,304$. Over 8 heads:
  $8 \times 98\,304 = 786\,432$.

  The heads are concatenated, giving $H \cdot D_V = 8 \times 64 = 512$ features, which the output
  projection maps back to $D$: $W_0 \in \mathbb{R}^{512 \times 512} = 262\,144$ entries.

  Total: $786\,432 + 262\,144 = 1\,048\,576$. Note that with $H \cdot D_V = D$, which is the usual
  choice, this is exactly $4D^2$ --- the cost of the block is quadratic in the model width and does
  not depend on the number of heads.
\end{solution}
\answerspace{4.0cm}

% =============================================================== Question 6 ===
\question
Once a Markov random field prior is added, segmentation can no longer be solved by deciding each
pixel independently. Please explain briefly why exact inference becomes hard, and describe one
approximate method presented in the lecture. You do not need to write down the equations.

\begin{solution}
  Answer: \emph{Why it is hard.} Without the MRF the posterior factorises over pixels, so each
  pixel's label can be chosen independently and the answer is found in one pass. The pairwise term
  couples neighbouring labels, so the posterior no longer factorises: the best label for a pixel
  depends on its neighbours' labels, which depend on theirs, and so on across the whole image.
  Computing the normalisation constant would require summing over every possible labelling of the
  image, of which there are exponentially many, and exact energy minimisation with categorical
  labels is NP-hard in 2D and above.

  \emph{An approximate method --- Gibbs sampling.} Sampling from the full posterior is infeasible,
  but sampling from the conditional distribution of a \emph{single} voxel given all the others is
  easy, because that conditional only involves the voxel's likelihood and its immediate
  neighbours, and normalising it means summing over the few possible labels of one voxel. So:
  initialise the labelling (for instance from the likelihood alone), then repeatedly visit each
  voxel and re-sample its label from that conditional. It is simple to implement and effective, but
  it is stochastic --- every run gives a different result --- the outcome depends on the
  initialisation and the parameters, and convergence can be slow.

  Iterated conditional modes and graph cuts are the other standard alternatives.
\end{solution}
\answerspace{7.5cm}

% =============================================================== Question 7 ===
\question
Which of the following statements on the differences between Active Shape Models (ASM) and Active
Appearance Models (AAM) is \textbf{inaccurate}?

\begin{choices}
  \choice ASM searches \emph{around} the current position, sampling along the surface normals of
  the model points, while AAM evaluates its fit \emph{at} the current position.
  \choice ASM minimises the distance between the model points and the corresponding points found
  in the image, whereas AAM minimises the difference between a synthesised model image and the
  target image.
  \choice AAM usually has a larger capture range than ASM, because the additional texture
  information gives it more to lock onto when the initialisation is poor.
  \choice AAM builds its texture model by warping every training image to the mean shape before
  sampling and normalising the intensities.
\end{choices}

\begin{solution}
  Answer: C.

  It is \textbf{ASM} that has the larger capture range, and (A) explains why: by sampling along the
  normals it actively looks a certain distance away from where the model currently sits, so it can
  find a boundary that is not underneath the model yet. AAM only compares its synthesised
  appearance with the image at the current position, so if the initialisation is too far off there
  is nothing to pull it in the right direction. ASM is also the faster of the two and locates
  feature points more accurately; what AAM buys in exchange is a better match to the image texture.
\end{solution}

% =============================================================== Question 8 ===
\question
Which of the following statements on non-linear transformation models is \textbf{incorrect}?

\begin{choices}
  \choice A naive parameterisation with one independent displacement vector per voxel can generate
  transformations that are physically implausible and even change the topology of the anatomy.
  \choice In a flow model each point follows the ordinary differential equation
  $dx/dt = v(x,t)$, and if the velocity field is smooth the resulting transformation is guaranteed
  to be a diffeomorphism.
  \choice A stationary velocity field has more parameters than a time-dependent one, which is why
  the scaling-and-squaring scheme is needed to make it tractable.
  \choice The diffusion model requires the displacement field to satisfy the Poisson equation,
  $\Delta u(x') + F(x') = 0$.
\end{choices}

\begin{solution}
  Answer: C.

  A stationary velocity field has \emph{fewer} parameters, not more: dropping the time dependence
  means storing one velocity vector per voxel instead of one per voxel \emph{per time point}. And
  scaling-and-squaring is not a workaround for a cost problem --- it is an efficient integration
  scheme that becomes available precisely \emph{because} the field is stationary. Both halves of
  the statement are inverted.

  The value of (B) is worth remembering separately: a diffeomorphism is smooth, invertible and has
  a smooth inverse, so a transformation built this way physically cannot tear or fold tissue. This
  is the guarantee that the naive parameterisation in (A) lacks, and it is the foundation of LDDMM.
\end{solution}

% =============================================================== Question 9 ===
\question
Which of the following statements on intensity normalisation is \textbf{incorrect}?

\begin{choices}
  \choice Min-max normalisation and matching the mean and standard deviation are both linear
  mappings of the intensities, and neither can change the contrast of an image.
  \choice Nyul's method matches the \emph{locations} of landmarks of the histograms, which aligns
  the distributions without making the two histograms identical.
  \choice Histogram matching exploits the spatial correlation between neighbouring pixels, which is
  what makes it robust to differences in the size of the imaged object.
  \choice Intensity variation is a particular problem for MRI because it has no absolute intensity
  scale, unlike CT where intensities are expressed in Hounsfield units.
\end{choices}

\begin{solution}
  Answer: C.

  Both halves are wrong. Histogram matching knows nothing at all about image content: it treats
  every pixel as an independent sample and uses \emph{no} spatial information or correlation
  between neighbours --- which is why matching image \emph{gradients} has been proposed as a
  remedy. And differences in object size are a genuine problem for it rather than something it is
  robust to: if one image contains much more of a given tissue than the other, a perfect histogram
  match is not even desirable, and mode matching is preferred.

  The other caveat worth knowing alongside this one is pathology: a large tumour changes the
  histogram, so matching a pathological image to a healthy one is hard, and robust statistics or
  outlier detection help.
\end{solution}

% ============================================================== Question 10 ===
\question
In patch-based label fusion the label at a target voxel is decided by
$L(x) = \arg\max_l \sum_i w_i\, \delta(\text{label}_i = l) \,/\, \sum_i w_i$, where the weights
$w_i$ come from the similarity between the patch around the target voxel and patches taken from
the atlas library. Which of the following statements about this scheme is correct?

\begin{choices}
  \choice The weights are computed from the spatial distance between the target voxel and the atlas
  voxel, not from the image content.
  \choice Only a coarse rigid or affine registration of the atlases is required, because
  correspondence is re-established locally by comparing patches inside a search window around each
  voxel.
  \choice The scheme requires exactly one atlas; with several atlases the weights no longer sum to
  one and the decision becomes ill-defined.
  \choice Because the weighting is borrowed from non-local means, the scheme can only be applied to
  denoising and not to the propagation of labels.
  \choice The search window has to cover the whole image, since restricting it would bias the
  weights.
\end{choices}

\begin{solution}
  Answer: B.

  This is the point of patch-based methods. In classical multi-atlas segmentation the
  correspondence between the atlas and the patient is whatever the registration says it is, so the
  result is only as good as a non-linear registration --- which is expensive and can fail. Patch
  matching replaces that: after a coarse alignment, each target voxel looks around a local search
  window for atlas patches that actually resemble its own neighbourhood, and takes their labels
  weighted by that resemblance. Correspondence is decided by image content at the point of use
  rather than by a global transformation computed in advance.

  (E) is the reverse of the truth: restricting the search window is standard practice, for cost, and
  it also encodes the sensible prior that the correct correspondence is somewhere nearby.
\end{solution}

% ============================================================== Question 11 ===
\question
A U-Net for segmentation takes $128 \times 128$ inputs with one channel. All convolutions are
$3 \times 3$ with ``same'' padding; each encoder level ends with $2 \times 2$ max pooling of stride
2, and each decoder level begins by upsampling by a factor of 2. The encoder levels produce 32, 64
and 128 channels respectively, and the bottleneck produces 256 channels. Each skip connection
concatenates the encoder features onto the upsampled decoder features at the matching level. Circle
True or False:

\begin{parts}
  \part \TFT\ \ The feature maps at the bottleneck have a spatial size of $16 \times 16$.

  \part \TFF\ \ Because all convolutions use ``same'' padding, no spatial information is lost in
  the encoder, and the skip connections are therefore redundant.

  \part \TFT\ \ After upsampling the bottleneck to $32 \times 32$ and concatenating the skip
  connection from the matching encoder level, the first decoder convolution at that level receives
  $256 + 128 = 384$ input channels.

  \part \TFF\ \ A $3 \times 3$ convolution mapping those 384 input channels to 128 output channels
  has $3 \cdot 3 \cdot 384 = 3\,456$ weights, excluding biases.
\end{parts}

\begin{solution}
  (a) TRUE --- three poolings, each halving: $128 \rightarrow 64 \rightarrow 32 \rightarrow 16$.
  The ``same'' convolutions do not change the size, so only the pooling does.

  (b) FALSE --- and this is the point of the architecture. It is the \emph{pooling}, not the
  convolutions, that destroys spatial detail: by the bottleneck the feature maps are $16 \times 16$
  and the fine boundary information is gone. The skip connections exist to carry the
  high-resolution features from the encoder across to the decoder, so the network can combine the
  large receptive field of the deep layers with the spatial precision of the shallow ones.

  (c) TRUE --- concatenation adds the channel counts: 256 upsampled from the bottleneck plus the
  128 arriving through the skip connection.

  (d) FALSE --- the count leaves out the output channels. A convolution kernel connecting every
  input channel to every output channel has $k_1 \cdot k_2 \cdot C_{in} \cdot C_{out}$ weights, so
  the correct number is $3 \cdot 3 \cdot 384 \cdot 128 = 442\,368$ --- 128 times the stated figure.
  Note also that this is where much of a U-Net's parameter budget goes: the concatenation makes the
  decoder convolutions wide.
\end{solution}

% ============================================================== Question 12 ===
\question
Which of the following statements on interpretability and its uses are correct (multiple answers
are possible)?

\begin{choices}
  \choice Saliency maps change over the course of training, and can be used as a kind of
  fingerprint of how a model responds to its inputs.
  \choice In interpretability-guided active learning, a sample is treated as informative when the
  model appears to struggle with it; easy samples tend to produce clean, focused saliency.
  \choice A model-agnostic method such as LIME explains a prediction by fitting a simpler,
  interpretable model that approximates the classifier in a neighbourhood of the input, weighted by
  a proximity measure.
  \choice Class-prototype maximisation returns the image from the training set that most strongly
  activates a given class.
  \choice Grad-CAM can be applied without any access to the internals of the model, since it only
  requires the class scores.
\end{choices}

\begin{solution}
  Answer: (A), (B) and (C) are correct.

  (D) misdescribes the method. Class-prototype maximisation performs gradient \emph{ascent on the
  input} to synthesise a pattern that maximises a class activation; the result is not a training
  image and indeed is not necessarily a realistic image at all. Regularising it with a generative
  prior is what makes the synthesised inputs look natural.

  (E) puts Grad-CAM on the wrong side of the taxonomy. It is a \emph{gradient}-based method, and it
  needs the gradients of the class score with respect to the convolutional feature maps --- so it
  requires access to the internals and is a white-box method. LIME is the model-agnostic,
  black-box one.
\end{solution}

\end{questions}

\end{document}
